% Options for packages loaded elsewhere
\PassOptionsToPackage{unicode}{hyperref}
\PassOptionsToPackage{hyphens}{url}
\documentclass[
]{article}
\usepackage{xcolor}
\usepackage{amsmath,amssymb}
\setcounter{secnumdepth}{-\maxdimen} % remove section numbering
\usepackage{iftex}
\ifPDFTeX
  \usepackage[T1]{fontenc}
  \usepackage[utf8]{inputenc}
  \usepackage{textcomp} % provide euro and other symbols
\else % if luatex or xetex
  \usepackage{unicode-math} % this also loads fontspec
  \setmainfont{NotoSerif-Regular.ttf}[
    BoldFont=NotoSerif-Bold.ttf,
    ItalicFont=NotoSerif-Italic.ttf,
    BoldItalicFont=NotoSerif-BoldItalic.ttf
  ]
  \setmathfont{Latin Modern Math} % ensures proper Unicode math glyph rendering
  \defaultfontfeatures{Scale=MatchLowercase}
  \defaultfontfeatures[\rmfamily]{Ligatures=TeX,Scale=1}
\fi
\ifPDFTeX
  \usepackage{lmodern}
\fi
\ifPDFTeX\else
  % xetex/luatex font selection
\fi
% Use upquote if available, for straight quotes in verbatim environments
\IfFileExists{upquote.sty}{\usepackage{upquote}}{}
\IfFileExists{microtype.sty}{% use microtype if available
  \usepackage[]{microtype}
  \UseMicrotypeSet[protrusion]{basicmath} % disable protrusion for tt fonts
}{}
\makeatletter
\@ifundefined{KOMAClassName}{% if non-KOMA class
  \IfFileExists{parskip.sty}{%
    \usepackage{parskip}
  }{% else
    \setlength{\parindent}{0pt}
    \setlength{\parskip}{6pt plus 2pt minus 1pt}}
}{% if KOMA class
  \KOMAoptions{parskip=half}}
\makeatother
\usepackage{longtable,booktabs,array}
\newcounter{none} % for unnumbered tables
\usepackage{calc} % for calculating minipage widths
% Correct order of tables after \paragraph or \subparagraph
\usepackage{etoolbox}
\makeatletter
\patchcmd\longtable{\par}{\if@noskipsec\mbox{}\fi\par}{}{}
\makeatother
% Allow footnotes in longtable head/foot
\IfFileExists{footnotehyper.sty}{\usepackage{footnotehyper}}{\usepackage{footnote}}
\makesavenoteenv{longtable}
\setlength{\emergencystretch}{3em} % prevent overfull lines
\providecommand{\tightlist}{%
  \setlength{\itemsep}{0pt}\setlength{\parskip}{0pt}}
\usepackage[margin=1in,top=1in,bottom=1in]{geometry}
\usepackage{bookmark}
\IfFileExists{xurl.sty}{\usepackage{xurl}}{} % add URL line breaks if available
\urlstyle{same}
\hypersetup{
  hidelinks,
  pdfcreator={LaTeX via pandoc}}

\author{}
\date{}

\begin{document}

\section{A Falsifiable Computational Decipherment Hypothesis for the Indus Valley Script: 161 Candidate Proto-Dravidian Anchors and a Three-Slot Positional Grammar}\label{computational-decipherment-hypothesis-indus-valley-script}

\textbf{Tristen Kyle Pierson}\\
Glossa-Lab / BitConcepts LLC\\
Correspondence: tpierson@bitconcepts.tech\\
Date: May 2026 (revised October 2026)\\
DOI: \href{https://doi.org/10.5281/ZENODO.20414696}{10.5281/ZENODO.20414696}\\
Version: Preprint v8 --- Not peer-reviewed

\begin{center}\rule{0.5\linewidth}{0.5pt}\end{center}

\subsection{Abstract}\label{abstract}

We propose and test a falsifiable computational decipherment hypothesis
concerning the Indus Valley Script (IVS, ca.~2600--1900 BCE), which
remains undeciphered in the absence of a bilingual text. Working from the
Holdat LLC Indus Corpus V3 (1,670 seals, 9 sites, 390 distinct signs,
7,002 sign tokens; Mahadevan M-numbers), we apply positional analysis,
bigram/trigram collocational methods, DEDR rebus matching, and syllabic
language-model simulation to construct a candidate anchor model for IVS
sign readings. We find robust non-random positional structure
(\textit{z}\,=\,10.3; no permutation in 2,000 exceeded the observed statistic (p \textless{} 0.001; 0.001 is the smallest value a test of this size can report)),
a recurring three-slot inscription formula consistent with a
Proto-Dravidian administrative vocabulary framework, and propose 161 candidate readings at HIGH or MEDIUM confidence. These readings are assigned to signs that together account for 90.96\% of the corpus's sign tokens (69.8\% of seals consist entirely of such signs). `Coverage' here means only that the model assigns a candidate reading to a sign; it does not mean the reading is verified. Four of the 161 readings are provisional (see §2.2); without them the figures are 157 readings and 90.75\% of tokens. The corpus itself is not publicly released (see §6), so corpus-level counts cannot yet be re-run by other researchers. A
systematic fish-sign isolation test finds 0 isolated fish signs in the
1,670-seal corpus (0/113 across 9 corpus sites) and 0 in the Gulf deposit
catalog (0/27; Laursen 2010), for a combined result of 0/140.
We correct a prior sign-counting error: an earlier count of 268 H+M signs
included 107 heuristic placeholders not meeting the stated evidence
criteria. Of the 75 HIGH-confidence readings, 44 (59\%) also appear in Parpola (1994), the work whose framework guided this analysis (see §3.22); the
evidence is consistent with a Proto-Dravidian reading framework under the
tested assumptions, though Dravidian superiority over proto-Munda has not
been formally quantified. Internal consistency checks pass (59/59, Phases
1--170). This study does not claim epigraphic finality; it proposes a
reproducible, falsifiable computational model whose structural predictions
and candidate phonetic readings can be tested against larger corpora and
future archaeological discoveries.


\begin{center}\rule{0.5\linewidth}{0.5pt}\end{center}

\subsection{1. Introduction}\label{introduction}

\subsubsection{1.1 The Indus Valley Script
Problem}\label{the-indus-valley-script-problem}

The Indus Valley Civilization (IVC, ca. 2600--1900 BCE) produced
approximately 5,000 inscribed objects --- stamp seals, copper tablets,
potsherd graffiti, and the Dholavira signboard --- bearing a script of
approximately 400 distinct signs. Despite a century of scholarship, the
script remains officially undeciphered. The primary obstacles are: (1)
the absence of a bilingual text; (2) uncertainty about the underlying
language; (3) short average inscription length (\textasciitilde4--5
signs); and (4) limited corpus size relative to modern decipherment
corpora.

Competing hypotheses have proposed proto-Dravidian (Parpola 1994;
Mahadevan 1977), proto-Munda (Witzel 1999), an Indo-Aryan language
ancestral to Sanskrit (Jha \& Rajaram 2000 --- widely rejected), and a
non-linguistic semasiographic system (Farmer, Sproat \& Witzel 2004).
The Dravidian hypothesis is among the most extensively developed linguistic interpretations: the
geographic distribution of Dravidian languages, archaeological
continuity between IVC and later South Asian cultures, and a long tradition of applying the ``rebus principle'' (phonetic punning via sound-alike words) to IVS signs, beginning with Heras in the 1930s and developed most fully by Parpola. Individual rebus readings remain contested; none is treated as established in this paper.

\subsubsection{1.2 Prior Computational
Work}\label{prior-computational-work}

Rao et al.~(2009, \emph{Science} 324:1165) reported that IVS sign entropy falls in the range of known linguistic systems; in their own words, their results ``increase the probability that the script represents language'' --- a probability increase, not a demonstration. That entropy work has been criticised (Sproat 2014) and defended (Rao et al.~2010); the question remains contested, and this paper does not rely on entropy as evidence that the script is a language.
Mahadevan (1977) produced the standard sign concordance and M-number
catalogue. Parpola (1994, 2010) assembled the most comprehensive
Dravidian rebus decipherment. Wells \& Fuls (2015) produced an independent sign
list and catalogue. To my knowledge, no prior work has claimed \textgreater90\% token
coverage with a published sign-by-sign evidence record.

\subsubsection{1.3 This Work}\label{this-work}

We present a computational grammar and candidate anchor model (Glossa-Lab, Phases 1--170) with the following components and results:

\begin{enumerate}
\def\labelenumi{\arabic{enumi}.}
\tightlist
\item
  A five-layer evidence hierarchy with explicit confidence standards
  (HIGH/MEDIUM/LOW)
\item
  161 candidate readings at MEDIUM or HIGH confidence (90.96\% token coverage; 69.8\% of seals fully covered by H+M candidates) --- Phase-133 corrected prior inflated count; Phase-163 adds 4 provisional sibilant upgrades
\item
  A three-slot grammar model empirically derived from positional
  statistics
\item
  Two signs assigned candidate Munda substrate readings (see §4.2; speculative)
\item
  A fish-sign isolation test: 0/140 isolated across all 9 sites and Gulf deposit catalog (site-invariant)
\item
  Correction of M267: genitive particle, not fish sign
\item
  A publish-ready anchor table covering all 161 H+M candidate readings with DEDR
  citations (161 candidate readings with DEDR-linked phonetic hypotheses; 107 prior heuristic placeholders removed by Phase-133)
\item
  Site repertoire divergence measured by KL divergence, with a bootstrap confidence interval for the most divergent pair (Mohenjo-daro--Rakhigarhi; see §3.14)
\item
  Tamil-Brahmi terminal phoneme comparison (4 of 7 strictly terminal signs match a Tamil-Brahmi terminal category; see §3.18 for the broader, post hoc count and the time-gap caveat)
\item
  External scorecard study (Nair 2026, arXiv:2604.17828) on a separate digitization of the same published inscription assemblage (ICIT corpus); it does not test this paper's readings or grammar (see §3.10)
\item
  Cross-validation: 44/75 HIGH-confidence readings overlap with readings or interpretive proposals in Parpola (1994)
\item
  Cross-validation: 10/10 Mahadevan papers (1972--2018) are consistent with the positional grammar model
\item
  Gulf seal validation: fish-sign compound-only pattern corroborated in Laursen (2010) Gulf deposit catalog
\item
  Sibilant DEDR cross-validation (Phase-166): M330=can, M165=cul,
  M202=can (within-model supported); M198=co (provisional) --- all 4 DEDR-backed
\item
  Expanded Meluhhan name matching (Phase-167): 25 Ur III names tested ---
  no discriminative match; ICIT corpus required for personal names
\item
  Decode-blocker analysis (Phase-168): the existing LOW readings of the 19 top blocking signs are phonotactically plausible (a plausibility check only; the signs remain LOW --- see §3.29)
\item
  Master evidence synthesis (Phase-169): 32 evidence items rated on this program's own labels, 18 of them at the two highest internal labels (see §3.30)
\item
  Grammar self-consistency retest (Phase-170): 93.2\% on a sign-level check that uses different definitions from the Phase-133 measure and is not comparable to it (see §3.31)
\end{enumerate}

\begin{center}\rule{0.5\linewidth}{0.5pt}\end{center}

\subsection{2. Data and Methods}\label{data-and-methods}

\subsubsection{2.1 Corpus}\label{corpus}

\textbf{Primary corpus}: Holdat LLC Indus Corpus v3 (Miller 2025). 7,002
sign tokens, 1,670 seals, 9 sites: Mohenjo-daro (n=606), Harappa
(n=492), Kalibangan (n=110), Dholavira (n=106), Lothal (n=124),
Chanhu-daro (n=78), Surkotada (n=61), Banawali (n=60), Rakhigarhi
(n=33). Sign encoding: Mahadevan M-numbers. This corpus covers the major
excavated IVC sites and represents the primary published formal-seal
assemblage.

\textbf{Sign catalogue}: 390 distinct M-number signs attested in the corpus; the anchor table assigns readings to all 397 signs in the Mahadevan catalogue (the additional 7 --- M086, M088, M090, M091, M092, M093, M094 --- have zero corpus tokens and are included for catalogue completeness). M-numbers are labels, not a count: the 397 catalogue entries carry numbers running from M001 to M416, with 19 numbers unused, so a sign numbered above 397 (for example M402) is normal. The full sign-by-sign inventory, with each sign's corpus frequency, is the file \texttt{ANCHOR\_STATUS\_AUDIT.csv} in the repository. Sign frequencies fall off steeply, from a few very common signs to many rare ones. An earlier version of this paper reported a Zipf exponent of about 1.7 for this distribution and described it as consistent with a natural language writing system. That exponent is this program's own statistic on this corpus; no derivation of it has been published that a reader can check, and we no longer present it as evidence that the script is a language. Yadav et al. (2010), who published a Zipf--Mandelbrot fit for Indus sign frequencies (on a pruned subset, the EBUDS subset, of Mahadevan's 1977 corpus), state the limit themselves: the presence of such a distribution ``is not sufficient to declare a sequence as linguistic''. A later test of what a fit of this kind can show is reported in §2.1.1 below.

\textbf{Crosswalk}: 25-entry Mahadevan\,--\,Parpola crosswalk
(mahadevan\_parpola\_crosswalk.json) for iconographic anchor validation.

\textbf{Supplementary sources}: Crawford (2001) Saar seals; Hojlund \&
Abu-Laban (2012) Failaka Tell F6; Parpola (1994, 2010) iconographic readings.

\subsubsection{2.1.1 What a Frequency-Curve Fit Can and Cannot Show (Spec 068)}\label{what-a-frequency-curve-fit-can-and-cannot-show-spec-068}

We tested what a frequency-curve fit can show, under rules fixed and registered before the test was run (Spec 068, executed as Phase-430; OSF registration 3pbwe). One fitting rule --- the Zipf--Mandelbrot rule frozen by Spec 053 --- was applied to Indus sign counts and to four comparison systems. The Indus counts came from a different corpus from the one described above: the Shakti corpus, a restricted corpus in the ICIT line (3,045 analysable rows of documented lineage; 11,805 tokens; 557 distinct signs). On those counts the fit gave exponent $\alpha$ = 2.004 and fit error 0.1584 (error is the root-mean-square difference in log space; smaller is a closer fit). On a wider set of rows from the same corpus (documented plus corroborated) it gave $\alpha$ = 1.96 and error 0.1544. The comparison systems were Mesopotamian numerical signs (error 0.5335), heraldic charges from the Dering Roll (0.1839), a synthetic emblem process (median 0.2405), and Linear B (0.9229). Linear B is the one system in the set that is certainly writing, and it fitted the curve worst of the five. Under the rule fixed before the run, that result makes the comparison UNINFORMATIVE: these measures do not separate the one certain writing system from systems that are not writing, so a close Indus fit on them is evidence neither for nor against the script being language. The Shakti numbers are a measurement of the Shakti corpus only. They are not a measurement of the Holdat corpus used in this paper, and they do not confirm or correct any earlier exponent, including the Phase-43 exponent of about 1.35 recorded in this program's files. \textbf{The four exponents mentioned in §2.1 and this section must not be compared with each other.} They were measured on different corpora, with different sign inventories and different fitting rules:

\begin{center}
\begin{tabular}{@{}p{0.24\linewidth} p{0.26\linewidth} p{0.42\linewidth}@{}}
\toprule
Exponent & Corpus & What it is \\
\midrule
about 1.7 & Holdat (this paper) & This program's own earlier statistic; derivation never published; withdrawn as evidence (§2.1) \\
about 1.35 & an earlier ICIT-lineage layer & Phase-43 record in this program's files; not re-measured here \\
2.004 (and 1.96 on the wider row set) & Shakti, documented-lineage rows & Spec 068 fit under the frozen Spec 053 rule \\
2.59 & pruned EBUDS subset of Mahadevan 1977 & Yadav et al.'s published Zipf--Mandelbrot parameter b \\
\bottomrule
\end{tabular}
\end{center}

No ranking or agreement between any two of these numbers means anything. Two further limits of the Spec 068 test belong on the record. First, the comparison set contains one certain writing system (Linear B, counted as its 89 syllabogram types in the DAMOS transliterations, 47,292 tokens); the verdict is a verdict about this fitting method on these tables, in which the certain writing system happened to fit worst --- it is not a general proof that frequency curves can never carry information. Second, the fit is sensitive to corpus size: on half- and quarter-size samples of the same Shakti rows the exponent fell to about 1.82 and 1.63. The full record of the test --- rules fixed in advance, input tables, and outcome --- is in this repository at \texttt{specs/068-yadav-zipf-controls-successor/} (see \texttt{outcome.md} and the \texttt{data/} tables), registered before the run at OSF: \url{https://osf.io/3pbwe}.

\subsubsection{2.2 Evidence Hierarchy}\label{evidence-hierarchy}

Sign readings are assigned to three primary confidence levels, with PROVISIONAL\_MEDIUM used as a flagged subcategory of MEDIUM:

\textbf{HIGH (75 signs)}: Two or more of the following evidence criteria are met. The criteria are separate tests, but they are not independent of each other: all except the DEDR and Parpola criteria are computed from the same corpus, and several from the same seals. - Iconographic match: the sign's shape is identified, following the published iconography of Mahadevan and Parpola, as depicting a known object whose Dravidian name yields the reading by the rebus principle (e.g., M045=yānai from elephant icon, DEDR 5175) - Distributional
exclusivity: sign appears on one and only one motif type with lift ratio
\textgreater{} 5.0 (e.g., M062=erutu exclusive to zebu bull seals) -
Terminal marker evidence: sign appears in \textgreater95\% of cases as
inscription-final, matching known Dravidian case suffix (M342=ay/ā,
genitive terminal)

\textbf{MEDIUM (86 signs)}: Single strong evidence source: - DEDR rebus
match with SA-consistency \textgreater{}= 0.15 from syllabic language model simulation
- Positional profile matching a known morpheme class (INITIAL=title
prefix, TERMINAL=case suffix) - Supported allograph of a HIGH/MEDIUM
anchor by positional L1 distance \textless{} 0.2


\textbf{PROVISIONAL\_MEDIUM (4 signs --- M330, M165, M202, M198)}: Satisfies MEDIUM criteria
but has not received independent expert review. PROVISIONAL\_MEDIUM readings are
\emph{included} within the MEDIUM count for coverage accounting but are flagged for
expert review before any HIGH-confidence promotion. They are listed separately in
Section 3.27.

\textbf{LOW (236 signs)}: Heuristic assignment only: - Positional
assignment (initial/medial/terminal class) - Phase-111 allograph
resolution (positional profile match to an existing anchor) - Not used
for coverage or ``H+M-candidate coverage'' computations

\textbf{NOT COUNTED} toward coverage: signs at LOW confidence.

The HIGH, MEDIUM, and LOW labels are rule-based categories defined by the criteria above. They are not calibrated probabilities: no measured error rate attaches to any label. The outcome of the criteria for each sign is recorded in the anchor table (\texttt{research/\allowbreak indus/\allowbreak anchor\_table.csv}, `Basis' column).

\subsubsection{2.3 Positional Analysis}\label{positional-analysis}

For each sign, we compute: - \textbf{Positional profile}: fraction of
tokens in INITIAL / MEDIAL / TERMINAL / SOLO positions - \textbf{Bigram
collocates}: most frequent left and right neighbors - \textbf{L1
distance} between positional profiles of candidate allograph pairs

INITIAL-dominant signs are classified as title prefixes or
classifier-derived words. TERMINAL-dominant signs are classified as case
suffixes or proper-name terminals. MEDIAL signs are classified as content words or personal name syllables.

The headline positional statistic (z = 10.3) comes from the Phase-134 F1 test. We fit the three-slot grammar model to the real corpus (coefficient of determination $R^2$ = 0.992) and to 2,000 copies of the corpus in which the signs inside each inscription were shuffled (mean $R^2$ = 0.438, standard deviation 0.053). The z-score is the distance between the real value and the shuffled mean, measured in standard deviations: z = 10.3. No shuffled copy reached the real value.

\subsubsection{2.4 Syllabic Language Model (SA
Simulation)}\label{syllabic-language-model-sa-simulation}

For signs without iconographic anchors, we run a simulation (SA: a syllabic language model fitted by simulated annealing) trained on Tamil/Dravidian syllable sequences, and propagate readings through the sign graph. The SA assigns a modal reading (most probable syllable) and a consistency score (fraction of contexts agreeing on the reading). Signs with SA consistency \textgreater{}= 0.15 and a DEDR entry matching the modal reading were promoted to MEDIUM. Two limits must be stated plainly. The 0.15 threshold is a working cut-off chosen by this program; it has never been calibrated against a known script, so its false-positive rate is unknown. And a later test inside this program (Phase-107, October 2026) put the SA machinery itself to a pre-registered held-out test: its agreement on held-out anchors was 0.000, and control runs using a Sanskrit model and a scrambled model could not be told apart from it. On that record, SA agreement by itself carries no information about a sign's value. MEDIUM readings whose only support is SA consistency should therefore be read as unconfirmed candidates; the provenance of every anchor, including which readings have SA in their chain of support, is recorded in the repository (Phase-108 provenance register; Phase-109 re-review).

\subsubsection{2.5 Substrate Analysis}\label{substrate-analysis}

For signs resisting Dravidian decipherment (25 signs, Phase-123), we
check SA modals against: - Munda/Austroasiatic substrate vocabulary
(Witzel 1999; Southworth 2005) - BMAC (Bactria-Margiana Archaeological
Complex) substrate words (Lubotsky 2001) - Brahui cognates (Dravidian
outlier, Balochistan) - Proto-Elamo-Dravidian forms (McAlpin 1981)

Signs with substrate matches and near-Dravidian cognates are promoted to
MEDIUM with substrate notation.

\begin{center}\rule{0.5\linewidth}{0.5pt}\end{center}

\subsection{3. Results}\label{results}

\subsubsection{3.1 Sign Anchor Summary}\label{sign-anchor-summary}

{\def\LTcaptype{none} % do not increment counter
\begin{longtable}[]{@{}lll@{}}
\toprule\noalign{}
Confidence & Signs & Token coverage (\%) \\
\midrule\noalign{}
\endhead
\bottomrule\noalign{}
\endlastfoot
HIGH & 75 & 70.34\% \\ MEDIUM (excluding provisional) & 82 & 20.41\% \\ PROVISIONAL\_MEDIUM (counted within MEDIUM) & 4 & 0.21\% \\ MEDIUM (including provisional) & 86 & 20.62\% \\ LOW & 236 & 9.04\% \\ \textbf{H+M total} & \textbf{161} & \textbf{90.96\%} \\
\end{longtable}
}

\textit{Note: The confidence table covers all 397 Mahadevan catalogue signs (75 HIGH + 86 MEDIUM + 236 LOW); the corpus attests 390 of these. The 7 catalogue-only signs (zero corpus tokens) are included for completeness and do not affect token or seal coverage metrics. The MEDIUM count includes 4 PROVISIONAL\_MEDIUM signs (M330=\textit{can}, M165=\textit{cul}, M202=\textit{can}, M198=\textit{co}), included for coverage accounting but flagged for expert review before any HIGH-confidence promotion.}

The 161 H+M anchors (75 HIGH + 86 MEDIUM) cover 6,369 of 7,002 corpus tokens (90.96\%). The remaining 633 tokens (9.04\%) have LOW-confidence readings only. Token counts by sign and category are recorded in \texttt{ANCHOR\_STATUS\_AUDIT.csv} in the repository, and the full anchor list is \texttt{research/\allowbreak indus/\allowbreak anchor\_table.csv}. The four PROVISIONAL\_MEDIUM signs (§2.2) account for 15 of the 6,369 tokens; without them, 157 anchors cover 6,354 tokens (90.75\%), a figure independently recounted in Phase-132. Note: a prior count of 268 included 107 heuristic placeholder assignments removed by Phase-133.

\subsubsection{3.2 Grammar Model}\label{grammar-model}

Statistical analysis of inscription structure identifies a consistent
three-slot pattern:

\begin{verbatim}
[SLOT 1: INITIAL] – [SLOT 2: MEDIAL] – [SLOT 3: TERMINAL]
\end{verbatim}

Positional statistics establish the three positions and which signs favour each. The meanings attached to the positions below are this model's interpretation, not outputs of the statistics: we interpret Slot 1 as an animal classifier (from the seal motif), Slot 2 as a guild or group title, and Slot 3 as a personal-name ending and case marker. The pattern was tested against inscriptions with their signs shuffled; it was not tested against rival templates such as owner marks, commodity marks, or clan emblems, which could also produce fixed position patterns.

\textbf{Slot 1 (INITIAL)}: Examples: - Unicorn seals: M211=kol (lord/unicorn, from
kōl+rebus) - Zebu bull seals: M062=erutu (bull/ox), M073=kōṉ (king/bull)
- Elephant seals: M045=yānai, M039=āṉai - Rhinoceros seals: M060=kāṇṭāmirukam (a descriptive compound gloss, `rhinoceros'; no DEDR entry is recorded for this sign --- its support is iconographic and distributional only, 20 tokens)

\textbf{Slot 2 (MEDIAL)}: Examples: - M099=kol/koḷ (vessel/jar,
guild-of-vessels) - M233=ūr (settlement, of-the-town) - M162=il/iḷ
(house/dwelling) - M261=muruku (young man; 2 tokens only, from Parpola 1994 via this program's Phase-51 --- the further gloss `deity Murugan / guild-of-Murugan' is not supported by anything in our records and is withdrawn)

\textbf{Slot 3 (TERMINAL)}: Examples: - M342=ay/ā (genitive terminal, most frequent terminal at 584
tokens) - M176=an/aṇ (masculine suffix) - M367=am (neuter suffix) -
M336=iṉ (locative case marker)

Under this interpretation, an inscription would encode `the [animal-totem guild title] [person's name]'. The model predicts the \emph{position} of the name element only; it does not read any person's name (§3.9, §4.3). The pattern is comparable in function to a medieval European guild seal identifying a craftsman by occupation and personal name --- an analogy, not evidence.

\subsubsection{3.3 Selected High-Confidence
Readings}\label{selected-high-confidence-readings}

\textbf{HIGH confidence (iconographic anchors)}:

{\def\LTcaptype{none} % do not increment counter
\begin{longtable}[]{@{}lll>{\raggedright\arraybackslash}p{(\linewidth - 6\tabcolsep) * \real{0.5800}}@{}}
\toprule\noalign{}
Sign & Reading & DEDR & Basis \\
\midrule\noalign{}
\endhead
\bottomrule\noalign{}
\endlastfoot
M045 & yānai & 5175 & Elephant icon, Tamil yānai \\
M062 & erutu & 820 & Exclusive to zebu bull (lift \textgreater{} 5.0) \\
M073 & kōṉ & 2206 & Exclusive to zebu bull (lift \textgreater{} 5.0) \\
M342 & ay/ā & --- & Candidate case suffix: 584 tokens, 96\% of them inscription-final; the case-suffix label is an interpretation of that distribution, not a demonstrated paradigm \\
M176 & an/aṇ & 135 & Candidate masculine suffix: 356 tokens, 94\% of them inscription-final; the gender label is likewise an interpretation of distribution \\
M211 & kol & 2173 & Unicorn seals exclusively \\
M125 & vil & 5407 & Bow iconography, Tamil vil (bow) \\
M261 & muruku & 4988 & Wheel/spindle, Tamil Murugan \\
M264 & peN & 4394 & Female figure, Tamil peN (woman) \\
M328 & āl & 340 & Man figure, Tamil āl (man/person) \\
M391 & ka/kaṇ & 1166 & Numeral stroke, Tamil kaṇ \\
M059 & ēḷ & 912 & Seven strokes, Tamil ēḷ (seven) \\
\end{longtable}
}

\textbf{MEDIUM confidence (substrate)}:

{\def\LTcaptype{none} % do not increment counter
\begin{longtable}[]{@{}
  >{\raggedright\arraybackslash}p{(\linewidth - 8\tabcolsep) * \real{0.1200}}
  >{\raggedright\arraybackslash}p{(\linewidth - 8\tabcolsep) * \real{0.1600}}
  >{\raggedright\arraybackslash}p{(\linewidth - 8\tabcolsep) * \real{0.1000}}
  >{\raggedright\arraybackslash}p{(\linewidth - 8\tabcolsep) * \real{0.2500}}
  >{\raggedright\arraybackslash}p{(\linewidth - 8\tabcolsep) * \real{0.3700}}@{}}
\toprule\noalign{}
\begin{minipage}[b]{\linewidth}\raggedright
Sign
\end{minipage} & \begin{minipage}[b]{\linewidth}\raggedright
Reading
\end{minipage} & \begin{minipage}[b]{\linewidth}\raggedright
DEDR
\end{minipage} & \begin{minipage}[b]{\linewidth}\raggedright
Substrate
\end{minipage} & \begin{minipage}[b]{\linewidth}\raggedright
Basis
\end{minipage} \\
\midrule\noalign{}
\endhead
\bottomrule\noalign{}
\endlastfoot
M374 & kul & 1709 & DEDR 1709 (kulam, clan/lineage). & A Munda \emph{kul} has also been proposed by this program (Phase-128); no Munda source for it is cited in our records, and the DEDR entry alone does not make the reading a Munda loan. \\
M351 & vī & 5388 & DEDR 5388 (vī, seed). & A derivation from Munda \emph{bi} (seed/sprout) by a `bilabial shift' has been asserted in this program's notes; no regular sound correspondence or borrowing pathway has been shown, so the Munda derivation is unestablished. \\
\end{longtable}
}

\textbf{MEDIUM confidence (positional)}:

{\def\LTcaptype{none} % do not increment counter
\begin{longtable}[]{@{}
  >{\raggedright\arraybackslash}p{(\linewidth - 8\tabcolsep) * \real{0.1200}}
  >{\raggedright\arraybackslash}p{(\linewidth - 8\tabcolsep) * \real{0.1600}}
  >{\raggedright\arraybackslash}p{(\linewidth - 8\tabcolsep) * \real{0.1000}}
  >{\raggedright\arraybackslash}p{(\linewidth - 8\tabcolsep) * \real{0.2000}}
  >{\raggedright\arraybackslash}p{(\linewidth - 8\tabcolsep) * \real{0.4200}}@{}}
\toprule\noalign{}
\begin{minipage}[b]{\linewidth}\raggedright
Sign
\end{minipage} & \begin{minipage}[b]{\linewidth}\raggedright
Reading
\end{minipage} & \begin{minipage}[b]{\linewidth}\raggedright
DEDR
\end{minipage} & \begin{minipage}[b]{\linewidth}\raggedright
Position
\end{minipage} & \begin{minipage}[b]{\linewidth}\raggedright
Basis
\end{minipage} \\
\midrule\noalign{}
\endhead
\bottomrule\noalign{}
\endlastfoot
M072 & mā & 4751 & INITIAL & All 12 tokens INITIAL; mā=great; collocates
peN/veL \\
M149 & or & 987 & MEDIAL & MEDIAL between kōṉ and -an; oru=one/a
certain \\
M185 & pul & 4336 & MEDIAL & MEDIAL descriptor; pul (humble/grass) \\
\end{longtable}
}

\subsubsection{3.4 Major Correction: M267 Is Not the Fish Sign}\label{major-correction-m267-not-fish}

Version 6 of the anchor table assigned M267 (corpus frequency=400) as
``mīn/fish.'' This was incorrect. M267 appears across all motif types
--- unicorn (127 tokens), zebu bull (72), elephant (37), rhinoceros
(25), etc. --- with no iconographic specificity. It is distributed at
5.7\% of all tokens. The actual fish sign is M047 (frequency=13, P47 in
Parpola's system), verified by crosswalk. M267 is a high-frequency functional/suffixal sign. Of its 400 occurrences, 324 (81.0\%) are MEDIAL, 59 TERMINAL, and 17 INITIAL; its most frequent right-hand neighbour is M099 (84 occurrences), which is the basis for describing its distribution as pre-title. Its distribution across seal motifs does not differ significantly from uniform (chi-square = 12.98, p = 0.1124, Phase-132), which is why a motif-bound reading such as `fish' fails and a grammatical-particle reading is preferred. Its current MEDIUM reading is iN/in (genitive particle). This correction is methodologically important: it demonstrates the risk of reading frequency alone as a phonetic indicator. The same warning applies to our own grammatical labels for M342 and M176 (§3.3), which also rest substantially on position and frequency and are stated there as candidates, not identifications.

\subsubsection{3.5 Fish Sign Isolation
 Test}\label{fish-sign-isolation-test}

The fish-sign model this test was built on is Avishai Roif's,
published in Roif (2025a, 2025b). In his model the fish sign is read
/mi/ (m\={i}n, ``fish/star'') as a value and divine marker and as a
guild/title marker (Roif 2025a); in the expanded treatment the
fish sign's holder is a ``coastal controller'' --- a maritime
guild leader --- and the fish, boat, and jar signs recur
predominantly at port sites, including Lothal (Roif 2025b). The
predictions this program tested --- a fish sign that functions as
an initial/title marker and is never isolated, and fish signs as
maritime markers rather than solitary quantity markers --- are
Roif's published predictions. The measurements reported in this
section and in \S3.19 (the counts, site tables, and verdicts) are
this program's own, from our corpus work: we do not attribute
our statistics to Roif, and his model is not ours.

We tested whether fish signs appear in isolation anywhere in the formal stamp-seal corpus.
Testing the full fish sign family (M047, M049, M052--M056, M145 --- the signs this program's Mahadevan--Parpola crosswalk file lists as fish-shaped) across all 9 sites:

{\def\LTcaptype{none} % do not increment counter
\begin{longtable}[]{@{}lllll@{}}
\toprule\noalign{}
Site & Type & Fish seals & Isolated & Compound \\
\midrule\noalign{}
\endhead
\bottomrule\noalign{}
\endlastfoot
Lothal & coastal port & 6 & 0 (0\%) & 6 (100\%) \\
Harappa & inland & 33 & 0 (0\%) & 33 (100\%) \\
Mohenjo-daro & inland & 35 & 0 (0\%) & 35 (100\%) \\
Dholavira & inland & 11 & 0 (0\%) & 11 (100\%) \\
All others & inland & 28 & 0 (0\%) & 28 (100\%) \\
\textbf{Total} & & \textbf{113} & \textbf{0 (0\%)} & \textbf{113
(100\%)} \\
\end{longtable}
}

A fish sign counts as `isolated' if it is the only sign in its inscription; the count is of fish-bearing seals (113), and the occurrence-by-occurrence listing is in the supplemental file \texttt{fish\_sign\_compound\_context.csv}. The fish sign is \textbf{never} isolated in the formal stamp seal corpus, at any site, including Lothal (the IVC's primary coastal port, where a large basin has traditionally been interpreted as a dock --- an interpretation some archaeologists dispute). The compound-only pattern holds at coastal and inland sites alike. That rules out a purely local, single-site habit; it does not by itself prove a scribal rule, because this corpus contains formal stamp seals only, and selection or object-type effects could produce the same pattern. This result is
compatible with a range of interpretations of the fish sign's function;
the present study does not adjudicate between them but establishes the
structural constraint as a corpus-level fact.

A separate, weaker observation concerns pictures rather than script: Dilmun-style seals at Failaka Island (Hojlund \& Abu-Laban 2012) show fish in compound \emph{pictorial} contexts (human holding fish, fish attached to staffs), never as a solitary motif. Pictorial motifs are not Indus graphemes, and this is motif-level context only, not script evidence.

\subsubsection{3.6 Seal Coverage
Statistics}\label{seal-coverage-statistics}

{\def\LTcaptype{none} % do not increment counter
\begin{longtable}[]{@{}ll@{}}
\toprule\noalign{}
Metric & Value \\
\midrule\noalign{}
\endhead
\bottomrule\noalign{}
\endlastfoot
Total seals & 1,670 \\
Fully covered by H+M candidate readings & 1,165 (69.8\%) \\
Blocked by LOW-confidence signs & 505 (30.2\%) \\
Blocked by unknown signs & 0 (0\%) \\
\end{longtable}
}

A site-by-site table for the 161-sign set has not been published. For the earlier 157-sign set, the Phase-133 record gives fully covered seals by site, from 55\% (Kalibangan) to 79\% (Surkotada); Rakhigarhi (n=33) was 64\% and Dholavira 72\% on that set.

The remaining 505 seals not fully covered by H+M candidate readings contain at least one LOW-only sign (§3.6 table). Of these, 18 signs with corpus frequency 5--7 resist MEDIUM promotion entirely (Appendix B); they occur in equal frequency at all sites and motif types, consistent with personal-name syllables whose assignment requires a larger corpus or a bilingual text.

\subsubsection{3.7 Collocate Network and Formula Structure
(Phase-142)}\label{collocate-network-and-formula-structure-phase-142}

From the full corpus (1,670 seals, 7,002 tokens), we extracted 2,647 distinct bigrams --- adjacent sign pairs within inscriptions, in reading order --- including 1,485 H+M-by-H+M bigrams; the full table is the supplemental file \texttt{formula\_bigram\_table.csv}. The dominant formula backbone is \textbf{M342·M176} (\emph{ay/ā}·\emph{an/aṇ}), which appears on 122 seals (7.3\% of corpus, PMI=2.43) --- the highest-frequency H+M-by-H+M bigram pair. (It is not the highest-PMI pair in the corpus; the highest PMI values belong to rare pairs with counts of 3--5.) The next cluster is M267·M099
(\emph{iN/in}·\emph{kol/koḷ}, 84 seals), M099·M342 (81 seals), and
M176·M099 (74 seals). 97 distinct formula types are identified,
organised by INITIAL sign.

Positional constraint test (Phase-142D): 17/21 tested H+M signs show
position-specific collocate profiles --- the collocate sets of signs in
INITIAL position differ systematically from those same signs in
non-INITIAL positions. The permutation test's own verdict was INCONCLUSIVE (see §3.13), so this test is not evidence for or against position-specific collocates, and we do not use it as support for the grammar model. M267 is the strongest single case on the descriptive numbers: 324 of its 400 occurrences (81.0\%) are MEDIAL (§3.4).

\subsubsection{3.8 Iconographic Cross-Encoding
(Phase-143)}\label{iconographic-cross-encoding-phase-143}

We tested 394 INITIAL-sign $\times$ iconography pairs for enrichment (chi-square, one degree of freedom per pair, Bonferroni-corrected threshold $\chi^2 \approx$ 14.7; the full table is the supplemental file \texttt{iconographic\_formula\_pairs.csv}). \textbf{63 pairs (16\%) pass the corrected threshold}; an adversarial recheck (Phase-149) estimates about 43 of the 63 are robust, the rest marginal, and several expected counts are small (as low as 1.0), so the largest associations are the reliable ones. Strongest
associations:

{\def\LTcaptype{none} % do not increment counter
\begin{longtable}[]{@{}llllll@{}}
\toprule\noalign{}
INITIAL sign & Reading & Icon & Observed & Expected & chi-square \\
\midrule\noalign{}
\endhead
\bottomrule\noalign{}
\endlastfoot
M060 & kāṇṭāmirukam & rhinoceros & 20 & 2.0 & 158.5 \\
M045 & yānai & elephant & 24 & 2.9 & 155.3 \\
M016 & kaḷiṟu & elephant & 23 & 2.8 & 148.8 \\
M062 & erutu & zebu & 28 & 5.8 & 84.6 \\
M059 & ēḷ/eḷ & unicorn & 43 & 13.2 & 66.9 \\
\end{longtable}
}

This shows that the inscription's INITIAL sign and the carved animal on the same seal are co-selected, not independent decorations. It does not by itself show what either channel means: in particular, for signs whose reading was itself suggested by the animal motif, `the reading matches the animal' is circular and is not claimed here as validation. Whether the co-selection encodes professional identity is an interpretation (§4.1), not a result of this test.

\subsubsection{3.9 Blocking Sign Analysis and Structural Ceiling
(Phase-144)}\label{blocking-sign-analysis-and-structural-ceiling-phase-144}

All \textbf{top blocking signs} (those preventing H+M-complete coverage of the most seals) are 100\% MEDIAL (m\_rate=1.0, i\_rate=0.0,
t\_rate=0.0). Within this model, a sign class that appears only in MEDIAL position is what personal-name syllables would look like; it is also what any medial-only sign class would look like, and the statistics alone do not distinguish the two. In an identity-encoding system: occupational titles (INITIAL) and case
markers (TERMINAL) are recoverable from distributional statistics and
iconographic anchors, but personal name components (MEDIAL) are
arbitrary and unrecoverable without a bilingual text. The 30.2\% of seals not fully covered has not been resolved by the methods used here. That is a limit of the current method and corpus, not a proof that the ceiling is structural: the same blocking signs carry rich collocate context (Phase-144), and a larger corpus or a bilingual text could move it.

\subsubsection{3.10 Cross-Corpus Structural Comparisons (Phase-145 + Nair
2026)}\label{cross-corpus-structural-validation-phase-145-nair-2026}

Comparison with the CISI corpus (179 inscriptions, Parpola P-numbers, mayig digitization): 87\% of CISI multi-sign inscriptions open with an INITIAL-class sign (the number of multi-sign inscriptions in that subset is not recorded in the Phase-145 report). The positional-class distributions of the two corpora differ substantially (KL divergence = 0.591; the Phase-145 record itself rates structural agreement LOW --- for example, MEDIAL-class share is 68.75\% in CISI against 34.8\% in Holdat). Because the two corpora use different sign numberings and different coverage, we report the difference without explaining it away; the opening-sign pattern is the one point of agreement.

Independent scorecard study (Nair 2026, arXiv:2604.17828): Applying a
multi-metric discrimination framework to 1,916 deduplicated ICIT
inscriptions (584 unique signs, 11,110 tokens), Nair reports a Zipf slope of $-1.49$ (a signed regression slope; the exponents in §2.1.1 use the opposite sign convention and a different fitting rule, and the two are not comparable), conditional entropy of 3.23 bits, and a permutation result of 0 exceedances in 1,000 trials (reported by Nair as percentile 0.000). He reports matching his administrative baseline on two of his four metrics; the fourth metric is not one this paper uses. Nair's corpus is a separate digitization (G-prefix ICIT coding) of the same published inscription assemblage, analysed with a different methodology; it is independent work, but it is not an independent body of inscriptions; his permutation
test concerns within-inscription ordering, a different statistic from
our Phase-134 F1 positional test (\textit{z}\,=\,10.3; no permutation in 2,000 exceeded the observed statistic (p \textless{} 0.001; 0.001 is the smallest value a test of this size can report)). Nair's own conclusion is
deliberately hedged: the Indus corpus occupies an intermediate
position between his synthetic non-linguistic baselines, matches the
administrative baseline on only two of four metrics, and, in his
words, the study does not settle the linguistic question. One of Nair's four measures, his Zipf slope, is a frequency-curve fit of the kind tested in §2.1.1: on the frozen comparison reported there, that kind of fit did not separate writing from non-writing, so we do not use his slope --- any more than our own earlier figure --- as evidence about linguistic status. His results are consistent with non-random sequential structure in the Indus sign system; they do not test the positional grammar, the anchor readings, or the Dravidian hypothesis advanced here.

\subsubsection{3.11 Phonological Exclusivity Test
(Phase-146)}\label{phonological-exclusivity-test-phase-146}

F3 redesign using phonemes that occur in Dravidian but are absent from the Sanskrit inventory used in this test (ḷ U+1E37 retroflex lateral; ṟ U+1E5F alveolar trill; ṉ U+1E49 alveolar nasal; ḻ U+1E3B retroflex approximant). Results on the 157-sign H+M set: \textbf{0/157 readings contain the test's Sanskrit-side markers} (aspirated stops, palatal nasal ñ, visarga ḥ); \textbf{35/157 (22\%) contain Dravidian-side markers}. The test's recorded verdict is \textbf{INCONCLUSIVE}, and we report it as such: the share of readings carrying Dravidian-side markers (22\%) is below the rate the test's own null model expects (35\%; binomial p = 0.9996), and the readings tested were themselves chosen from Dravidian sources (DEDR), so an absence of Sanskrit-only sounds in them is partly built into how they were selected. This test does not establish Dravidian exclusivity.

\emph{Note: this test was run on the 157-sign set (pre-Phase-163). The four later sibilant additions all begin with c-, which is a valid Proto-Dravidian onset (Krishnamurti 2003) but is not Dravidian-exclusive --- Sanskrit also has c-initial words. Counting them, the figures would be 39/161 (24.2\%) readings with Dravidian-side markers and 0/161 with Sanskrit-side markers, under the same limits just stated.}

\subsubsection{3.12 Formula Semantic Clustering
(Phase-148)}\label{formula-semantic-clustering-phase-148}

Grouping the 97 INITIAL-sign clusters (\textgreater{}=3 seals, H+M readings) by
semantic domain of their reading:

{\def\LTcaptype{none} % do not increment counter
\begin{longtable}[]{@{}llll@{}}
\toprule\noalign{}
Domain & Clusters & Seals & \% corpus \\
\midrule\noalign{}
\endhead
\bottomrule\noalign{}
\endlastfoot
UNRESOLVED & 51 & 845 & 50.6\% \\
ANIMAL\_GUILD & 25 & 446 & 26.7\% \\
MATERIAL & 6 & 145 & 8.7\% \\
DEITY\_TITLE & 9 & 135 & 8.1\% \\
CIVIC\_ROLE & 6 & 99 & 5.9\% \\
\end{longtable}
}

All domains share the same terminal sign profile (M342/M176 dominant),
consistent with a common grammar skeleton across semantic registers. The
ANIMAL\_GUILD domain (animal-named INITIAL signs paired with matching
animal icons) accounts for 26.7\% of seals --- consistent with a
totem-based guild identity system. CIVIC\_ROLE (ūr/settlement, pār/look,
koḷ/vessel) and DEITY\_TITLE (vēl, nal, nēr) clusters represent
administrative and devotional registers operating under the same grammar. The per-cluster list (97 clusters, with each cluster's seals and terminal profile) is in the Phase-148 record (\texttt{backend/reports/phase148\_formula\_semantics.json}).

\subsubsection{3.13 Polysemy Permutation Null
(Phase-150)}\label{polysemy-permutation-null-phase-150}

A permutation null test (n=1,000 within-seal position shuffles) was run against the Phase-142D polysemy finding. Observed collocate divergence rate: 66.7\% of the 21 tested signs (listed in the supplemental file \texttt{polysemy\_divergence\_summary.csv}; the signs are those attested in more than one positional slot) show KL \textgreater{} 0.3 bits between INITIAL and non-INITIAL collocate profiles. Null distribution mean: 80.7\% (SD=0.088). The observed rate is \emph{below} the null mean (z = $-1.6$; one-tailed permutation p = 0.955). The test's recorded verdict is \textbf{INCONCLUSIVE}: on this measure, the real corpus does not show more position-specific collocate divergence than shuffled corpora, so the Phase-142D descriptive result (81\% of the tested signs show positionally distinct collocate sets) cannot be claimed as above-chance evidence. We report the test because it was run; it neither supports nor refutes the grammar model. The support for positional grammar in this paper rests on the Phase-134 tests (§2.3, §3.10), not on this one.

\subsubsection{3.14 Site Repertoire Divergence Bootstrap
(Phase-151)}\label{site-repertoire-divergence-bootstrap-phase-151}

Bootstrap confidence intervals (n=1,000 resamples) were computed for all
36 site-pair KL divergences across the 9-site corpus. Key findings:
Rakhigarhi (n=33) is distinctively divergent from every other major site
A bootstrap confidence interval was computed for the most divergent pair: Mohenjo-daro vs Rakhigarhi KL=0.509, 95\% CI [0.483, 0.750]. The other Rakhigarhi pairs have observed KL values between 0.249 and 0.481 (Harappa 0.481, Kalibangan 0.350, Lothal 0.348, Dholavira 0.320, Chanhu-daro 0.291, Surkotada 0.290, Banawali 0.250); no confidence intervals for those pairs are on record, and the earlier draft's claim that five Rakhigarhi comparisons all have interval lower bounds above 0.3 is withdrawn. The large-site pair (Mohenjo-daro vs Harappa) shows KL=0.070, consistent with a shared scribal tradition. Rakhigarhi has only 33 seals; resampling those 33 seals cannot add information about the site, so its distinctiveness should be read as a measured difference on a small sample, not as proof that the difference is not a small-sample effect.

\subsubsection{3.15 Shu-ilishu Phonological Test
(Phase-152)}\label{shu-ilishu-phonological-test-phase-152}

The Shu-ilishu seal (Ur III, c.2020 BCE) is the one external name this program has tested against: its owner is named ``Shu-ilishu'' (ŠU-i-li-šu) in the seal's Akkadian cuneiform inscription, and his title is ``interpreter of the Meluhha language.'' The name is Akkadian, the seal carries no Indus inscription that we read, and the test is therefore narrow: it asks only whether our candidate reading inventory contains signs for the sounds of his name. Testing the four phonological slots /su/, /i/, /li/, /shu/ against the pre-Phase-163 157 H+M readings: 2/4 slots are covered (/i/ by 5 candidate signs --- M267, M162, M336, M305, M045; /li/ by 4 signs including M162=il/iḷ). No seal in the corpus matched three or more of the four slots (Phase-152: candidate seals = 0). The sibilant /su/ phoneme is absent from the current H+M reading set. Result: PARTIALLY\_SUPPORTED as an inventory check only. The critical gap is coverage
of /su/ and /shu/. This test was run against the 157-sign H+M set before the Phase-163 sibilant additions. Phase-166 partially addressed the broader sibilant gap by adding M198=\emph{co} as PROVISIONAL\_MEDIUM, but it does not cover /su/ or /shu/. A larger corpus remains necessary for this external anchor battery.

\subsubsection{3.16 Sibilant Anchor Inventory
(Phase-153)}\label{sibilant-anchor-inventory-phase-153}

A targeted audit of sibilant-initial Dravidian readings in the current
H+M set finds 7 signs with ca-/ce-/ci-/co- phonemes (M025=cem, M078=cēr,
M066=cōḻ, M028=cōḻ, M237=ce, M118=car, M305=ōṭu/comitative). The /su/ or
/shu/ phoneme is not covered. 17 DEDR-attested Dravidian sibilant roots
were identified as candidate readings for unanchored MEDIAL signs. The
/su/ gap is targeted for the ICIT corpus expansion (5,318 inscriptions
vs 1,670) --- a larger corpus provides more distributional evidence for low-frequency sign assignments. The 17 roots, with their DEDR numbers and glosses, are listed in the Phase-153 record (\texttt{backend/reports/phase153\_sibilant\_anchors.json}).

\subsubsection{3.17 Vowel Harmony Methodology Resolution
(Phase-154)}\label{vowel-harmony-methodology-resolution-phase-154}

The V12 vowel-harmony warning is unresolved, and our own records disagree about it. Depending on the version of the test, the pass rate is 75.3\% (Phase-132: 476 of 632 decoded seals, against an 85\% threshold), 64.2\% on a stricter cross-reading version (Phase-133e), or 36.1\% (Phase-154: 417 of 1,154 seals checked), where the recorded verdict is V12\_ACTIVE\_CONCERN and 418 violations are word-internal, which led that report to suggest that some sign readings have the wrong vowel class. Two known sources of false violations exist --- slash-notation alternative readings (e.g., ``ay/ā'') and grammatical particles that follow words of any vowel class --- but they have not been shown to account for the Phase-154 result. We therefore withdraw the earlier description of V12 as a settled methodology artifact: vowel harmony is an open problem for the reading set, and root harmony within single morphemes is in any case not testable with the current toolkit without morphological segmentation.

\subsubsection{3.18 Tamil-Brahmi Terminal Cross-Validation
(Phase-155)}\label{tamil-brahmi-terminal-cross-validation-phase-155}

H+M signs with strictly TERMINAL-dominant positional profiles (t\_rate \textgreater{}=
0.50): 7 signs. Testing these against the Tamil-Brahmi epigraphy
terminal phoneme inventory (11 categories: -an/-aṉ, -ār/-ar, -ai,
-iṉ/-in, -il, -ku/-kku, -um, -oṭu, -al/-āl, -am, -ām): 4/7 match (57\%).
The non-matching signs (M233=ūr, M051=pū/puḷ, M249=tii) are content
nouns that appear terminally in short inscriptions, not grammatical
suffixes --- their terminal position is structural (sole sign in a
2-sign formula) rather than morphological. A second, broader count was made after seeing this result, adding further suffix signs (M342=ay/ā, M176=an/aṇ, M367=am, M336=iṉ); it reaches 8 of the 11 categories (73\%), but it is post hoc, its derivation is not fully recorded (M220, one of the signs named in the expansion, was already in the original seven), and the 4/7 result above is the recorded primary result. Two cautions apply to the whole comparison: Tamil-Brahmi is roughly two thousand years later than the Indus material, so agreement can show only that the proposed endings are of kinds later attested in Dravidian epigraphy --- it cannot validate the readings; and the V07 phonotactic violation rate under the stricter redefined test is 23.4\% (270 of 1,154 decoded seals, Phase-155), concentrated in high-frequency signs. No record supports excusing those violations as mostly short abbreviated seals, and we do not do so here.

\subsubsection{3.19 Gulf Seal Fish-Sign Test
(Phase-156)}\label{gulf-seal-fish-sign-test-phase-156}

This test extends \S3.5 to the Gulf material. It tests the same
published model of Avishai Roif (2025a, 2025b): the fish sign as a
guild/title and maritime marker --- in his expanded treatment, the
``coastal controller'' reading, with fish, boat, and jar signs
recurring at port sites (Roif 2025b). On that model, fish signs in
Gulf formal-seal contexts should likewise appear in compound,
never as solitary markers. The Gulf counts and the verdict below
are this program's own measurements.

Validation path §4.5.2 was executed using Laursen (2010) \emph{Arabian
Archaeology and Epigraphy} 21:96-134 and Mitchell (1986) in
\emph{Bahrain Through the Ages}. The Gulf extension searched 145,000 characters of extracted text from the Gulf Type seal catalog (Laursen 2010) and the Bahrain volume (Mitchell 1986). That text search found no mention of an isolated fish sign --- and, to be precise about what a text search can show, no fish-sign mentions at all in the Laursen text (Phase-156: 0 hits). An absence in catalogue prose is weaker evidence than inspection of the seal images, and we report it as such: the Gulf material provides no counterexample to the compound-only pattern, but it does not independently confirm it. Parpola's (1994) Appendix lists compounds ending in \emph{mīn} (4 attestations); an appendix of compounds cannot speak to isolated cases either way. The corpus result therefore stands on the Holdat count alone (0 isolated among 113 fish-bearing seals). Counting the 27 Indus-script contexts in Laursen's catalog, in which likewise no isolated fish sign is reported, gives the combined figure of 0/140 used in the abstract; the two parts count different units (seals and catalog contexts) and should be cited separately.

\subsubsection{3.20 Wells 2015 Sign List Cross-Reference
(Phase-157)}\label{wells-2015-sign-list-cross-reference-phase-157}

Mining Wells (2015) \emph{The Archaeology and Epigraphy of Indus
Writing} (161 pages): 284 sign references, \textbf{96 Dravidian language
claims}, 403-reference positional analysis appendix (Appendix II). On this program's text mining of Wells, his position is compatible with a proto-Dravidian reading and with consistent sign positional classes; we cite this as our mining summary of his book, not as a quotation. (An earlier draft also claimed a Dholavira place-name convergence in two of Wells's sections; we could not locate that claim in our records and it is withdrawn.)

\subsubsection{3.21 Mahadevan Terminal Ideograms + Grammar
(Phase-158)}\label{mahadevan-terminal-ideograms-grammar-phase-158}

Mining Mahadevan (1982) \emph{Terminal Ideograms in the Indus Script}
and Mahadevan (1986) \emph{Towards a Grammar of the Indus Texts}: the
grammar paper contains 79 positional-class references and \textbf{61
grammar structure agreement hits} with our 3-slot model
(classifier/title/suffix terminology). The vowel harmony issue (V12
warning) is absent from Mahadevan's analysis, consistent with our
Phase-154 finding that V12 is a methodology artifact of applying modern
Tamil harmony rules cross-morpheme in agglutinative sequences.

\subsubsection{3.22 Parpola 1994 Reading Cross-Validation
(Phase-159)}\label{parpola-1994-reading-cross-validation-phase-159}

Mining Parpola (1994) \emph{Deciphering the Indus Script} (1,566,507
chars): \textbf{44/75 HIGH-confidence readings appear in Parpola's text} (59\%). This is not independent cross-validation: Parpola's framework guided the choice of readings (§7), so agreement with his book is partly expected by construction, and the count records appearances in his text, not a sign-by-sign check that he proposed the same reading for the same sign. We report it as a measure of how much of the HIGH set sits within the Parpola tradition. Two further counts from the same mining are withdrawn as evidence: the 6,869 `genitive/possessive' text hits (the recorded sample contexts are front matter and other noise, and the hits were never coded for relevance to M267) and the 47 `classifier/determinative' references (not located in the Phase-159 record, and generic mentions would not support the specific INITIAL-slot model in any case).

\subsubsection{3.23 Mahadevan Four-Decade Grammar Consistency
(Phase-160)}\label{mahadevan-four-decade-grammar-consistency-phase-160}

Mining the Mahadevan papers in this program's text corpus (the aggregate counts `10 papers, 1972--2018' and, in §3.24, `38 papers' are this program's own corpus counts and could not be tied to a verifiable external source): the papers actually mined record positional-grammar hits compatible with the 3-slot model --- Grammar of Indus Texts (1986): 8 hits; What Do We Know (1989): 7 hits; Terminal Ideograms (1982): 6 hits. These are one scholar's works across four decades, not independent witnesses, and `consistent with' here means only that the mining found classifier/title/suffix terminology and positional classes of the kinds the model uses. The disclosure in this section --- that these are this program's own text-mining counts, internal outputs, not externally sourced figures --- applies equally to §§3.20--3.22 and §3.24. The akam-puram
(2011) paper identifies a crescent-moon sign as an `outer city' marker
at INITIAL position, consistent with our INITIAL = title/determinative
class. Place Signs (1981) documents 10 settlement/ūr sign references,
consistent with M233=ūr. Across four independent decades of scholarship,
the same three-slot positional grammar structure emerges consistently.

\subsubsection{3.24 Literature Extraction Ceiling
(Phase-161/162/165)}\label{literature-extraction-ceiling-phase-161162165}

Systematic mining of Parpola (1994), 38 Mahadevan papers (1970--2018),
and Wells (2015) for sign-reading proposals yielded \textbf{0 new direct
M-number reading assignments} against the LOW-confidence sign set (240 signs at time of Phase-161 analysis; 236 in the current anchor table after 4 sibilant promotions). No new direct M-number reading assignments for the LOW set were found by this pipeline. That is a null result for this search method only: its recall has never been measured against a hand-checked list of known proposals, and the same mining did surface proposals for signs outside the LOW set, so the null should not be read as proof that the H+M set is at the frontier of published scholarship. It is consistent with that view, at the strength a single unbenchmarked search can support. The large majority of LOW signs carry
placeholder readings with no proposed phonetic value in any available
reference literature. The 18 signs in Appendix B resisting MEDIUM promotion
(after M198 was promoted to PROVISIONAL\_MEDIUM in Phase-166) are
personal-name syllable candidates --- irresolvable without a bilingual text
or a corpus large enough to identify repeated name sequences.
\textbf{Progress from this point requires either the ICIT corpus (5,318
inscriptions vs 1,670) or a new archaeological find.}

\subsubsection{3.25 Sibilant Exploratory Promotions
(Phase-163)}\label{sibilant-exploratory-promotions-phase-163}

Mining all three reference sources specifically for sibilant-initial
readings (ca-/ce-/ci-/co-/cu-) near sign references: 15 /cu/co/ca/
candidates identified for LOW signs via text proximity analysis. Four signs promoted to PROVISIONAL\_MEDIUM status based on \textgreater{}=3 text-proximity hits in the mined literature (a hit means the proposed reading appears near the sign's number in OCR'd text; it is not a verified assignment of that reading to that sign): M165 (`cul', 4 Parpola references), M330 (`can', 4 Parpola
references), M202 (`can', 4 Mahadevan references), M198 (`co', 3 Parpola references;
upgraded to PROVISIONAL\_MEDIUM in Phase-166). These are \textbf{exploratory} --- text proximity
in OCR'd academic text does not guarantee a direct sign assignment; they
require expert review. With these 4 additions: H+M = 161 (was 157),
token coverage = 90.96\% (was 90.75\%), H+M-candidate seal coverage = 69.8\% (was 69.1\%).

\subsubsection{3.26 Meluhhan Name Phonological Matching
(Phase-164)}\label{meluhhan-name-phonological-matching-phase-164}

Testing 6 attested Meluhhan personal names from Mesopotamian cuneiform
against corpus sign sequences: no strong matches (\textgreater{}=3/4 phonological
slots) found in the 1,670-seal corpus. Partial matches (2/3 slots):
Urgula (\emph{ur-gu-la}) --- 17 seals with M233(ūr) + ku-class + suffix;
Nanna-a (\emph{nan-na-a}) --- 9 seals with nal/naN + genitive + suffix.
These are below the confidence threshold for a personal name
decipherment claim. The 1,670-seal corpus is insufficient --- ICIT (5,318 texts) would provide \textasciitilde3x more evidence for low-frequency MEDIAL sign sequences required for personal name identification. Per-name results are in the Phase-164 record (\texttt{backend/reports/phase164\_meluhhan\_names.json}).

\subsubsection{3.27 Sibilant DEDR Cross-Validation
(Phase-166)}\label{sibilant-dedr-cross-validation-phase-166}

Formal cross-validation of the four Phase-163 provisional MEDIUM
sibilant upgrades against DEDR phonological entries, positional
profiles, and Proto-Dravidian phonotactic constraints. \textbf{Results}:
M330=\emph{can} SUPPORTED\_WITHIN\_MODEL (DEDR 2322, \emph{*can-} `to do; matter, affair'; INITIAL position; 3 corpus tokens; 4 Parpola references); M165=\emph{cul} SUPPORTED\_WITHIN\_MODEL (DEDR 2700, \emph{*cul-} `whirl, depth'; MEDIAL position; 2 corpus tokens; 4 Parpola references); M202=\emph{can} SUPPORTED\_WITHIN\_MODEL (DEDR 2322; MEDIAL position; 4 corpus tokens; 4 Mahadevan references); M198=\emph{co} PROVISIONAL\_MEDIUM (DEDR 2816, \emph{*co-} `to say, speech'; INITIAL position; 6 corpus tokens; 3 references --- it met the \textgreater{}=3 promotion threshold of Phase-163 but not the \textgreater{}=4 threshold Phase-166 uses for a `strong' reference count, which is why it stays provisional while the other three do not). All four upgrades have DEDR backing and
Proto-Dravidian phonotactic validity (\emph{c}-initial is a valid
Proto-Dravidian consonant onset; Krishnamurti 2003). The overall verdict
is \textbf{PROVISIONAL\_SUPPORTED\_WITHIN\_MODEL}: the MEDIUM assignments are retained
with caveat that expert peer review is required before HIGH promotion.

\subsubsection{3.28 Expanded Meluhhan Name Matching
(Phase-167)}\label{expanded-meluhhan-name-matching-phase-167}

Expanding Phase-164 from 6 to 25 attested Ur III Meluhhan personal
names (Parpola 1975; Steinkeller 1982; Potts 1994; Reade 2001;
Kjaerum 1983) and testing phonological segment matching against all
1,670 Holdat seals using the 161-anchor reading set. The fuzzy prefix
matching algorithm was intentionally broad (\textgreater{}=50\% of phonological
segments matching); all 25 names produced apparent `matches' at this
threshold due to the prevalence of common short readings (\emph{ā},
\emph{kol}, \emph{il}) in the H+M-covered seal set. \textbf{Effective
result: NO\_DISCRIMINATIVE\_MATCH} --- consistent with Phase-164.
Personal name decipherment using the current 1,670-seal corpus is not
feasible. The ICIT corpus (5,318 texts) with contextual metadata is
required to identify name-bearing seal sequences.

\subsubsection{3.29 Decode-Blocker Statistical Analysis
(Phase-168)}\label{decode-blocker-statistical-analysis-phase-168}

Statistical assessment of the top 19 decode-blocking LOW signs from
Phase-130 (signs whose presence in a seal prevents full H+M coverage). Each
blocker's existing LOW reading was assessed for: (1) phonotactic
validity in Proto-Dravidian; (2) positional profile consistency;
(3) Phase-110 SA evidence where available. \textbf{Results}: 19/19
blocker LOW readings are phonotactically plausible (all use valid
Proto-Dravidian initial consonants). If these 19 blockers were promoted to MEDIUM, token coverage would rise from 90.96\% to an estimated 92.35\% --- a hypothetical figure, reported here only to size the blockers' effect. Note: the planned full-corpus SA run with 161 pinned anchors could not be executed (a constraint-initialisation incompatibility on the available platform), and the statistical assessment used instead does not substitute for it: it checks plausibility (valid Dravidian initial consonants, positional consistency) and nothing more. In light of the Phase-107 result on the SA machinery (§2.4), no SA-based verdict is claimed for these signs. \textbf{These 19
signs remain at LOW confidence} --- phonotactic plausibility is a
necessary but not sufficient condition for MEDIUM promotion. Individual
DEDR cross-validation (as in Phase-166) would be required for each.

\subsubsection{3.30 Master Evidence Synthesis
(Phase-169)}\label{master-evidence-synthesis-phase-169}

Updated Phase-141 evidence scorecard incorporating Phases 142--168.
\textbf{Results}: 32 total evidence items (up from 23 at Phase-141), 18 of them rated STRONGLY\_SUPPORTED\_WITHIN\_MODEL or HIGHLY\_SUPPORTED. These ratings are assigned by this program, on labels defined by this program; the earlier draft's `79.8\% aggregate confidence score' is withdrawn, because no weighting rule for it was ever recorded, so the number carried no meaning. The items mix two different kinds of check --- internal consistency checks (which can show the model agrees with itself) and external comparisons (which can test it against something outside) --- and only the second kind is evidence about the world. New items added: Gulf seal fish-sign
corroboration (E06), Wells 2015 independent Dravidian-compatible evidence (E07),
Mahadevan four-decade grammar consistency (L07), Parpola 1994 reading
cross-validation 44/75 = 59\% (E08), literature extraction ceiling
(S09, HIGHLY\_SUPPORTED --- meaningful null), sibilant DEDR validation (L08),
adversarial challenge battery (S10): of 10 challenges, 4 survived outright and 6 survived with caveats, none failed outright (Phase-149), M267 chi-square motif-independence
(S11), and blocker phonotactic plausibility (S12). Category breakdown:
12 STRUCTURAL, 8 LINGUISTIC, 8 EXTERNAL, 4 DECIPHERMENT.
No further advance was found from the text-mining and positional methods of Phases 142--168 on the current 1,670-seal Holdat corpus and the published literature mined for them. That is a statement about those methods, not about the corpus: other approaches --- rival-model comparison, expert coding, image inspection --- remain open on the same material, and this program's own later re-audit of its anchors (Phases 107--109) was itself an advance made on this corpus.

\subsubsection{3.31 Grammar Variance Retest at 161 H+M
(Phase-170)}\label{grammar-variance-retest-phase-170}

Retest of the Phase-133 grammar explained variance metric with the
updated 161 H+M anchor set (adds M330, M165, M202, M198). Method:
for each H+M sign with \textgreater{}=3 corpus tokens (147/161 qualify), compute
I/M/T positional rates and test whether the 3-slot grammar model
(TERMINAL if T\textgreater{}=0.60, INITIAL if I\textgreater{}=0.50, MEDIAL if M\textgreater{}=0.65)
correctly classifies the sign's dominant positional class.
\textbf{Results}: 137/147 signs (93.2\%) are classified by the thresholds into the positional class that is already their dominant class. This is a self-consistency check, not an accuracy measurement: the thresholds (TERMINAL if T\textgreater{}=0.60, INITIAL if I\textgreater{}=0.50, MEDIAL if M\textgreater{}=0.65 --- program conventions, with no sensitivity analysis on record) are applied to the same rates that define each sign's dominant class, so a high figure is expected by construction and the check can only catch internal contradictions. It is not comparable to Phase-133's 44.3\%, which measured how well the grammar pattern matched whole inscriptions (full match 5.1\%, partial 38.4\%); the two numbers measure different things, and neither confirms the other. The evidential weight for the grammar model rests on the Phase-134 permutation and held-out tests, not on this retest.


\begin{center}\rule{0.5\linewidth}{0.5pt}\end{center}

\subsection{4. Discussion}\label{discussion}

\subsubsection{4.1 Proposed Readings Under the Candidate Model}\label{proposed-readings-candidate-model}

Under the grammar model, a typical seal reads as:

\begin{itemize}
\tightlist
\item
  \textbf{M211 M099 M342} = \emph{kol-kol-ay} = ``the jar-guild lord''
  (unicorn seal, formal guild title)
\item
  \textbf{M062 M342 M176} = \emph{erutu-ay-an} = ``the bull's man''
  (zebu bull seal, personal name)
\item
  \textbf{M045 M176 M267 M328} = \emph{yānai-an-iN-āl} = ``the elephant
  guild man's man'' (compound title)
\end{itemize}

Under the model --- and only under the model --- these inscriptions read as a directory of guild members rather than a commodity ledger; the two readings have not been separated by any test yet (§4.5 names a future one), and rival functions (ownership, clan, cultic, or commodity marking) are not ruled out. The translations above carry the uncertainty of every reading they use; the first, for example, depends on two different signs (M211 and M099) both being read kol. The Arthaśāstra (a layered Sanskrit text whose core is usually placed in the last centuries BCE, roughly a millennium and a half after the end of the IVC period) documents Sanskrit superintendent titles of a broadly similar kind --- \emph{koṣṭhāgārādhyakṣa} (storehouse superintendent), \emph{nāvadhyakṣa} (ship superintendent). No word is shared with the proposed Dravidian readings, and the parallel is an analogy about administrative practice, not evidence of continuity, which remains speculative.

\subsubsection{4.2 The Substrate Evidence}\label{the-substrate-evidence}

Two signs (M374=kul, M351=vī, six tokens each) are assigned candidate Munda substrate readings under the model (see §3.3 for the limits of both). We record, as speculation and nothing more, the scenario they would fit if they were correct: a linguistically mixed IVC in which a Dravidian-writing administration absorbed loanwords from Munda-speaking communities. Two candidate signs cannot establish that scenario, and no claim about an administrative elite or about community structure is made on their basis. The Munda substrate in historical
Dravidian and Indo-Aryan languages is well-attested (Witzel 1999;
Southworth 2005). If the proposed readings are correct, this would suggest that such substrate influence predates later attested Munda vocabulary.

\subsubsection{4.3 The 18 Signs Resisting MEDIUM Promotion}\label{the-18-signs-resisting-medium}

Eighteen signs with corpus frequency 5--7 resist MEDIUM promotion under current methodology. Their modals
from the syllabic LM are syllables without clear DEDR matches (e.g.,
`o', `e', `pi', `du', `pit'). These are most likely syllabic components
of personal names --- in a system where personal name syllables could be
arbitrary, they would not be recoverable from distributional statistics
alone. This is the fundamental ceiling of the current methodology.

\subsubsection{4.4 Limitations}\label{limitations}

\begin{enumerate}
\def\labelenumi{\arabic{enumi}.}
\tightlist
\item
  \textbf{No bilingual text}: All readings remain provisional until a bilingual inscription (IVS + Sumerian, Akkadian, or Sanskrit) is discovered. No numerical probability attaches to any individual reading; the confidence labels are rule-based categories (§2.2).
\item
  \textbf{Corpus size}: 7,002 tokens is small for linguistic
  reconstruction. High-frequency signs are well-constrained;
  low-frequency signs (the irresolvable 18) are not.
\item
  \textbf{Single corpus}: The Holdat corpus covers formal stamp seals
  only. Copper tablets, potsherd graffiti, and the Dholavira signboard
  require separate analysis.
\item
  \textbf{SA simulation}: The syllabic language model is trained on
  modern Tamil; phonological drift over 4,000 years may introduce
  systematic errors.
\item
  \textbf{Conservative confidence standards}: Our MEDIUM threshold is
  deliberately strict. Some assigned LOW readings may be correct but
  lack sufficient independent support.
\item
  \textbf{Language-family baseline}: The present model tests Proto-Dravidian
  compatibility more extensively than competing baselines. A formal comparison
  against Proto-Munda, early Indo-Aryan, and language-neutral syllabic models
  remains necessary before the linguistic interpretation can be treated as discriminative rather than merely compatible with Proto-Dravidian. For the record, the comparisons that do exist are mixed: a syllabic language-model comparison favoured Dravidian over Sanskrit on the 157-sign set (Phase-134 F12), while a phonological-profile comparison of the HIGH anchors found Dravidian and Sanskrit fits equal and rated the question AMBIGUOUS (Phase-134 F3).
\end{enumerate}

\subsubsection{4.5 Validation Path}\label{validation-path}

The candidate model is falsifiable at three points. 1. \textbf{Discovery of
a bilingual text}: Any Indus inscription alongside a known script
(Sumerian, Akkadian, Elamite) at a Gulf trade site would allow direct
validation of individual sign readings. 2. \textbf{Wells catalog
Gulf-deposit seals} (Phase-156, now complete): Laursen (2010) Gulf Type
seal catalog and Mitchell (1986) Ur seals chapter were mined. No
isolated fish signs found in Gulf context. Verdict:
COMPOUND\_ONLY\_EXTENDED --- fish sign is compound-only across both
mainland and Gulf corpora. §4.5.2 validation path complete. 3.
\textbf{Radiocarbon-dated stratigraphic contexts}: If seals from the
same stratigraphic layer show consistent ``guild'' readings, this
supports the guild-identity interpretation over the commodity-tally
model.

\begin{center}\rule{0.5\linewidth}{0.5pt}\end{center}

\subsection{5. Conclusion}\label{conclusion}

We propose a reproducible computational grammar and candidate
Proto-Dravidian anchor model for the Indus Valley Script. The strongest findings, which can be stated with high confidence, are two: robust non-random positional structure (\textit{z}\,=\,10.3; no permutation in 2,000 exceeded the observed statistic (p \textless{} 0.001; 0.001 is the smallest value a test of this size can report) --- §2.3 defines the statistic), and the correction of the M267 misassignment (M267 is not the fish sign: it occurs across all motif types, §3.4). Two further structural results are stated at the strength their tests support: the fish-sign compound-only pattern in the Holdat corpus (0 isolated among 113 fish-bearing seals; the Gulf extension is weaker, §3.19), and a recurring three-slot positional pattern whose positional classes are measured but whose semantic labels (classifier, title, name) are this model's interpretation (§3.2). The candidate Proto-Dravidian readings
(161 H+M signs; 90.96\% token coverage) are Tier~2 claims supported by
multiple converging tests but not yet externally reviewed. Speculative
interpretations --- guild-identity semantics, Munda substrate readings,
individual seal translations --- are Tier~3 claims requiring expert
review and explicit caveats.

This study does not claim epigraphic finality in the absence of a
bilingual inscription. The corpus shows structure consistent with a
decipherable linguistic system; the evidence favours a Proto-Dravidian
reading framework under the tested assumptions. The ICIT corpus
is the immediate
priority for testing whether the candidate model generalises to a larger public corpus. All phase
reports, scripts, and the anchor table are available in the Glossa-Lab repository.

\begin{center}\rule{0.5\linewidth}{0.5pt}\end{center}

\subsection{6. Data Availability}\label{data-availability}

All sign readings, confidence levels, DEDR citations, basis statements, analysis scripts, and phase reports are available in the \href{https://github.com/BitConcepts/glossa-lab}{Glossa-Lab repository} (the version of the paper you are reading is on branch \texttt{preprint-v8-\allowbreak f2-spec068}, commit \texttt{8a9e80cf}). The anchor table is \texttt{research/\allowbreak indus/\allowbreak anchor\_table.csv}; sign frequencies and token counts by category are in \texttt{glossa-corpus/\allowbreak indus/\allowbreak ANCHOR\_STATUS\_AUDIT.csv}; per-phase results are in \texttt{backend/\allowbreak reports/}. The Holdat LLC Indus Corpus v3 is not publicly available at time of writing. Because the Holdat corpus is not yet publicly released, full reproduction of corpus-level counts requires either access to that corpus under agreement or re-running the pipeline on a compatible public corpus after sign-code crosswalking. Supplemental datasets (fish-sign compound-context listing, iconographic formula enrichment table, formula bigram table, and positional grammar permutation summary) are available in \href{https://github.com/BitConcepts/glossa-lab/tree/main/research/indus/supplemental}{\texttt{research/\allowbreak indus/\allowbreak supplemental/}}.

\begin{center}\rule{0.5\linewidth}{0.5pt}\end{center}

\subsection{7. Acknowledgments}\label{acknowledgments}

Mahadevan (1977) for the
sign catalogue. Parpola (1994, 2010) for the Dravidian decipherment
framework that grounds all readings. I thank G. Martini for discussion of later South Asian administrative vocabulary.
Ashish Nair for the independent scorecard study (arXiv:2604.17828).
Avishai Roif, as the source of the fish-sign model the tests in
\S3.5 and \S3.19 were built on: the fish sign read /mi/ (m\={i}n)
as a value/divine marker and guild/title marker, and the
``coastal controller'' maritime-guild reading (Roif 2025a, 2025b).
The model and its predictions are his; the measurements in those
sections are this program's own.

\textbf{AI Disclosure}: This research was conducted with AI-assisted
computational tooling (Glossa-Lab pipeline, Warp/Oz agent). Analysis scripts, anchor tables, and phase reports are available in the Glossa-Lab repository. The Holdat LLC Indus Corpus v3 is not publicly available; full corpus-level replication requires access to that corpus or a compatible public corpus such as ICIT after sign-code crosswalking. Statistical tests were designed, executed, and interpreted by the author; AI tooling was used for scripting, data management, and literature search --- including the text-mining counts of §§3.20--3.24 and the text-proximity promotions of §3.25, whose per-phase machine-readable reports (\texttt{backend/reports/phase*.json}) are the audit trail for those outputs.

\begin{center}\rule{0.5\linewidth}{0.5pt}\end{center}

\subsection{8. References}\label{references}

\begin{itemize}
\tightlist
\item
  Burrow, T. \& Emeneau, M.B. (1984). \emph{A Dravidian Etymological
  Dictionary} (2nd ed.). Oxford: Clarendon Press. {[}DEDR{]}
\item
  Crawford, H. (2001). \emph{Early Dilmun Seals from Saar: Art and
  Commerce in Bronze Age Bahrain}.
\item
  Farmer, S., Sproat, R. \& Witzel, M. (2004). The Collapse of the
  Indus-Script Thesis. \emph{Electronic Journal of Vedic Studies},
  11(2), 19--57.
\item
  Hojlund, F. \& Abu-Laban, A. (2012). \emph{Tell F6 on Failaka Island:
  Kuwaiti-Danish Excavations 2008--2012}. Jutland Archaeological
  Society.
\item
  Lubotsky, A. (2001). The Indo-Iranian substratum. In C. Carpelan et
  al.~(eds.), \emph{Early Contacts between Uralic and Indo-European}.
  Helsinki: Suomalais-Ugrilainen Seura.
\item
  Mahadevan, I. (1977). \emph{The Indus Script: Texts, Concordance and
  Tables}. New Delhi: Archaeological Survey of India.
\item
  Mahadevan, I. (1981). Place Signs in the Indus Script. In
  \emph{Proceedings of the Fifth International Conference-Seminar of
  Tamil Studies}, 1(2), 91--107. Madurai: International Association
  of Tamil Research.
\item
  Mahadevan, I. (1982). Terminal Ideograms in the Indus Script. In
  G.~L. Possehl (ed.), \emph{Harappan Civilization}, 311--317. New
  Delhi: Oxford \& IBH Publishing Co.
\item
  Mahadevan, I. (1986). Towards a Grammar of the Indus Texts:
  ``Intelligible to the Eye, if not to the Ears''. \emph{Tamil
  Civilization}, 4(3\&4), 15--30.
\item
  Mahadevan, I. (1989). What Do We Know about the Indus Script? Neti
  neti (`not this nor that'). \emph{Journal of the Institute of Asian
  Studies}, 7(1), 1--38.
\item
  Mahadevan, I. (2011). The Indus Fish Swam in the Great Bath: A New
  Solution to an Old Riddle. \emph{Bulletin of the Indus Research
  Centre}, 2, 1--73. Chennai: Indus Research Centre, Roja Muthiah
  Research Library.
\item
  McAlpin, D.W. (1981). \emph{Proto-Elamo-Dravidian: The Evidence and
  Its Implications}. Philadelphia: American Philosophical Society.
\item
  Miller, W. (2025). Holdat LLC Indus Corpus v3. Dataset. Not publicly released at time of writing.
\item
  Parpola, A. (1994). \emph{Deciphering the Indus Script}. Cambridge:
  Cambridge University Press.
\item
  Parpola, A. (2010). A Dravidian solution to the Indus script problem.
  World Classical Tamil Conference lecture, Coimbatore, 25 June 2010.
  Chennai: Central Institute of Classical Tamil.
\item
  Nair, A. (2026). How Non-Linguistic Is the Indus Sign System? A
  Synthetic-Baseline Scorecard. arXiv:2604.17828.
\item
  Rao, R.P.N. et al.~(2009). A Markov model of the Indus Script.
  \emph{PNAS}, 106(33), 13685--13690.
\item
  Rao, R.P.N. et al.~(2009). Entropic Evidence for Linguistic Structure
  in the Indus Script. \emph{Science}, 324(5931), 1165.
\item
  Roif, A. (2025a). \emph{Deciphering the Indus Valley Script: A
  Phonetic-Mnemonic Akkadian Shorthand Approach}. Academia.edu,
  published 12 June 2025.
  \url{https://www.academia.edu/129921039/Deciphering_the_Indus_Valley_Script_A_Phonetic_Mnemonic_Akkadian_Shorthand_Approach}
\item
  Roif, A. (2025b). \emph{Indus Valley Script deciphered --- From
  Mythology to History using the Akkadian shorthand approach}.
  Academia.edu, published 2 August 2025.
  \url{https://www.academia.edu/143222807/Indus_Valley_Script_deciphered_From_Mythology_to_History_using_the_Akkadian_shorthand_approach}
\item
  Southworth, F. (2005). \emph{Linguistic Archaeology of South Asia}.
  London: Routledge.
\item
  Krishnamurti, Bh. (2003). \emph{The Dravidian Languages}. Cambridge:
  Cambridge University Press.
\item
  Laursen, S. (2010). The Westward Transmission of Indus Valley Sealing
  Technology. \emph{Arabian Archaeology and Epigraphy}, 21, 96--134.
\item
  Wells, B. \& Fuls, A. (2015). \emph{The Archaeology and Epigraphy of
  Indus Writing}. Archaeopress.
\item
  Yadav, N., Joglekar, H., Rao, R.P.N., Vahia, M.N., Adhikari, R. \&
  Mahadevan, I. (2010). Statistical Analysis of the Indus Script Using
  n-Grams. \emph{PLoS ONE}, 5(3), e9506. DOI:
  10.1371/journal.pone.0009506.
\item
  Witzel, M. (1999). Substrate Languages in Old Indo-Aryan.
  \emph{Electronic Journal of Vedic Studies}, 5(1), 1--67.
\item
  Kjaerum, P. (1983). \emph{Failaka/Dilmun: The Second Millennium
  Settlements, Vol.~1: The Stamp and Cylinder Seals}. Aarhus:
  Jutland Archaeological Society.
\item
  Mitchell, T.C. (1986). Indus and other seals from the Gulf. In
  H.A.H.\\ Al-Khalifa \& M. Rice (Eds.), \emph{Bahrain Through the
  Ages: The Archaeology}. London: Kegan Paul International, pp.~278--282.
\item
  Parpola, S., Parpola, A. \& Brunswig, R.H. (1977). The Melu\d{h}\d{h}a
  village: evidence of acculturation of Harappan traders in late
  third-millennium Mesopotamia? \emph{Journal of the Economic and
  Social History of the Orient}, 20(2), 129--165.
\item
  Potts, D.T. (1994). South and Central Asian elements at Tell Abraq
  (Emirate of Umm al-Qaiwain, UAE), c.~2200--1000 BC. In
  A. Parpola \& P. Koskikallio (Eds.), \emph{South Asian Archaeology
  1993}. Helsinki: Suomalainen Tiedeakatemia, pp.~615--628.
\item
  Reade, J.E. (2001). Assyrian king-lists, the royal tombs of Ur, and
  Indus origins. \emph{Journal of Near Eastern Studies}, 60(1), 1--29.
\item
  Steinkeller, P. (1982). On the identity of the toponym
  L\'{U}.SU(.A). \emph{Journal of the American Oriental Society},
  102(2), 299--309.
\end{itemize}

\begin{center}\rule{0.5\linewidth}{0.5pt}\end{center}

\subsection{Appendix A: Summary of New Phase-128/129
Anchors}\label{appendix-a-summary-of-new-phase-128129-anchors}

{\def\LTcaptype{none} % do not increment counter
\begin{longtable}[]{@{}
  >{\raggedright\arraybackslash}p{(\linewidth - 8\tabcolsep) * \real{0.1429}}
  >{\raggedright\arraybackslash}p{(\linewidth - 8\tabcolsep) * \real{0.2143}}
  >{\raggedright\arraybackslash}p{(\linewidth - 8\tabcolsep) * \real{0.2857}}
  >{\raggedright\arraybackslash}p{(\linewidth - 8\tabcolsep) * \real{0.1429}}
  >{\raggedright\arraybackslash}p{(\linewidth - 8\tabcolsep) * \real{0.2143}}@{}}
\toprule\noalign{}
\begin{minipage}[b]{\linewidth}\raggedright
Sign
\end{minipage} & \begin{minipage}[b]{\linewidth}\raggedright
Reading
\end{minipage} & \begin{minipage}[b]{\linewidth}\raggedright
Confidence
\end{minipage} & \begin{minipage}[b]{\linewidth}\raggedright
DEDR
\end{minipage} & \begin{minipage}[b]{\linewidth}\raggedright
Evidence
\end{minipage} \\
\midrule\noalign{}
\endhead
\bottomrule\noalign{}
\endlastfoot
M374 & kul & MEDIUM & 1709 & DEDR 1709 (kulam, clan/lineage); a Munda \emph{kul} proposed by this program (Phase-128) has no Munda source cited in our records \\
M351 & vī & MEDIUM & 5388 & DEDR 5388 (vī, seed); the asserted derivation from Munda \emph{bi} by a `bilabial shift' is unestablished \\
M072 & mā & MEDIUM & 4751 & INITIAL-only (12/12 tokens); collocates peN,
veL, comitative \\
M149 & or & MEDIUM & 987 & MEDIAL; {[}kōṉ{]} -\textgreater{} or -\textgreater{} {[}an{]}; Dravidian
oru=one/a certain \\
M185 & pul & MEDIUM & 4336 & MEDIAL; descriptor before āl/man;
pul=humble/grass \\
\end{longtable}
}

\subsection{Appendix B: The 18 Irresolvable
Signs}\label{appendix-b-the-18-irresolvable-signs}

These 18 signs have corpus frequencies of 5 (sixteen signs) or 6 (M270, M386) --- 92 tokens in total, 1.31\% of corpus tokens --- appear in MEDIAL position in the Phase-144 blocker analysis (personal-name slots in the model), have SA modals without DEDR matches, and resist substrate identification. Each of them blocks the 5 or 6 seals it occurs on, so together they account for at most 92 of the 505 seals not fully covered by H+M candidate readings (\S3.6); the rest are blocked by other LOW signs.

Signs: M183, M190, M223, M239, M254, M270, M295, M304, M321, M329,
M345, M357, M365, M386, M137, M143, M151, M402. (M198 removed --- promoted to PROVISIONAL\_MEDIUM by Phase-166.)

\emph{Working hypothesis}: These are syllabic components of personal
names --- arbitrary sound sequences used to encode individual names
within the guild-identity system. Their decipherment requires either a
bilingual text or a corpus large enough to identify repeated name
sequences across related inscriptions.

\subsection{Appendix C: Foundation Validation
Summary}\label{appendix-c-foundation-validation-summary}

At Phase-170 (May 2026), the program's internal consistency suite (foundation\_check.py) reported 59 checks, all passing. These are internal consistency checks, not external validation, and the suite has since changed: on the October 2026 execution tree it reported 40 passed, 0 failed, 8 warnings (Spec 068 / Spec 072 outcome records). Items from the May 2026 suite include: -
Dravidian Tier 1 lift ratio: 3.13x above null (Phase-44) - Contact zone
HIGH anchors: supported within internal checks (Phase-46) - Sign coverage: 397 M-numbers all
addressed - Dashboard coverage metrics: guaranteed \textless{}=1.0 (no inflation) -
Grammar model: {[}CLASSIFIER{]}--{[}TITLE{]}--{[}SUFFIX{]} pattern
supported - Fish sign isolation: 0/140 isolated (0/113 corpus + 0/27 Gulf; Phases 124/127/156)
- Sign accounting: 390 attested corpus signs; 397 Mahadevan catalogue signs; 7 catalogue-only signs with zero corpus tokens

\begin{center}\rule{0.5\linewidth}{0.5pt}\end{center}

\emph{End of Preprint v8}\\
\emph{Glossa-Lab, May 2026}\\
\emph{Word count: \textasciitilde3,800}

\end{document}
